\scanpage{17}
d'où aussitôt la flèche envisagée.

Il faudrait encore expliciter comment (7) (8) résultent de la donnée de l'opération de $\mathcal E'$ sur $\widetilde X'(J)$, induisant une opération \uline{libre} de $\pi_1\subset\mathcal E'$ sur $\widetilde X'(J)$ (ayant $J$ comme ens quotient)\ldots.

\uline{Reformulation}

Considérons les systèmes
\[
\left\{\begin{gathered}
\diagram{\begin{tikzcd}[ampersand replacement=\&]
1 \arrow[r,dash] \& \pi \arrow[r,dash] \& \mathcal E \arrow[r,dash] \& \mathcal G \arrow[r,dash] \& 1 \& \text{ext. de groupes}
\end{tikzcd}}\\
Z\quad\mathcal E\text{-ens }\inserted{\ne\varnothing},\text{ où }\pi\subset\mathcal E\text{ opère librement}
\end{gathered}\right.
\]
Ils forment une \uline{catégorie} $C$.
Considérons les systèmes $(\Pi,\widetilde\Pi,p,\mathcal G,a)$ où $\Pi$ est un groupoïde \inserted{connexe}, $p:\widetilde\Pi\longrightarrow\Pi$ un rev. universel de $\Pi$, et $\mathcal G$ un groupe opérant \deleted{\unsure{sur}} sur $\Pi$ par automorphismes par $a:\mathcal G\longrightarrow\operatorname{Aut}(\Pi)$. Ils forment une catégorie $C'$. On va définir des équivalences quasi-inverses l'une de l'autre
\diagram{\begin{tikzcd}[ampersand replacement=\&]
C \arrow[r,shift left,"\alpha"] \& C' \arrow[l,shift left,"\beta"]
\end{tikzcd}}
\mspage{16}
